\documentclass[letterpaper,11pt,twoside]{article}
\usepackage[utf8]{inputenc}
\usepackage{amsmath,amsfonts,amssymb,amsthm,latexsym}
\usepackage[spanish,es-noshorthands]{babel}
\usepackage[T1]{fontenc}
\usepackage{lmodern}
\usepackage{graphicx,hyperref}
\usepackage{tikz,pgf}
\usepackage{multicol}
\usepackage{marvosym}
\usepackage{fancyhdr}
\usepackage[height=9.5in,width=7in]{geometry}
\usepackage{fancyhdr}
\pagestyle{fancy}
\fancyhead[LE]{Colegio Arborizadora Baja}
\fancyhead[RE]{PEI:``Hacia una cultura para el desarrollo sostenible''}
\fancyfoot[RO]{\Email iedabgerman@autistici.org}
\fancyhead[LO]{\url{www.autistici.org/mathgerman}}
\fancyfoot[RE]{\Email cedarborizadoraba19@redp.edu.co}
\fancyfoot[LE]{Calle 59I N$^{\circ}$ 44A - 02 \Telefon 7313994 - 7313995}

\author{Germ\'an Avenda\~no Ram\'irez~\thanks{Lic. Mat. U.D., M.Sc. U.N.}}
\title{\begin{minipage}{.2\textwidth}
\includegraphics[height=1.75cm]{Images/logo-colegio.png}\end{minipage}
\begin{minipage}{.55\textwidth}
\begin{center}
Animaplano 3\\
Matemáticas $11^{\circ}$
\end{center}
\end{minipage}\hfill
\begin{minipage}{.2\textwidth}
\includegraphics[height=1.75cm]{Images/logo-sed.png} 
\end{minipage}}
\date{}
\thispagestyle{plain}
\begin{document}
\maketitle
Nombre: \hrulefill Curso: \underline{\hspace*{44pt}} Fecha: \underline{\hspace*{2.5cm}}
\begin{multicols}{2}
\section*{Cuestionario}
\begin{enumerate}
\item Los estudiantes de grado 11 han organizado una fiesta y por cada dos hombres que pagan entrada, ingresan gratis tres mujeres. Si en la fiesta hay 120 personas, ¿cuántas mujeres han entrado gratis?
\item En una empresa editorial 35 impresoras producen 205 copias en 10 minutos. ¿Cuántas copias producen en el mismo tiempo 14 impresoras de la misma marca y referencia?
\item En un fábrica de vehículos 9 empleados que trabajan 4 horas diarias terminan una tarea en 62 días. ¿Cuántos días demoran 3 empleados con el mismo rendimiento, en hacer la misma tarea trabajando 8 horas diarias?
\item Si se alimentan 7 gallinas mediante 2 raciones diarias, un bulto de alimento solo dura 7 días. ¿Para cuántos días dura un bulto si solo se alimenta una gallina con una ración diaria?
\item En un restaurante saben que cuando hay 267 clientes se necesitan 60 empleados en turnos de 6 horas cada uno. Manteniendo estas proporciones, ¿a cuántos clientes se puede atender si solo hay 15 empleados en turnos de 8 horas?
\item En una fábrica de helados se sabe que 16 neveras congelan 42 conos en 21 minutos. Manteniendo estas proporciones, ¿en cuántos minutos aumenta el proceso si se desea congelar 210 helados usando 28 neveras?
\item El valor del límite $\displaystyle{\lim_{x\rightarrow 3}x^{3}+x^{2}+12}$ es:
\item El valor de $\displaystyle{\lim_{x\rightarrow 7}x^{2}}$ es:
\item El valor de $\displaystyle{\lim_{x\rightarrow 7}x^{2}-9}$
\item El valor de $\displaystyle{\lim_{x\rightarrow 6}x^{2}+2}$
\item El séptimo número primo
\item El valor de $\displaystyle{\lim_{x\rightarrow 5}}\dfrac{x^{2}+3x-40}{x-5}$ es
\item La quinta potencia del primer número primo
\item El décimo tercer número primo
\item El valor de \quad $\displaystyle{(9)\cdot\lim_{x\rightarrow 5}
\left(\dfrac{x^{2}-2x-15}{x-5}\right)}$ es
\item Si $\displaystyle{\lim_{x\rightarrow a}f(x)}=9$ y $\displaystyle{\lim_{x\rightarrow a}g(x)}=7$, entonces $\displaystyle{\lim_{x\rightarrow a}f(x)\cdot g(x)}=$?
\item El décimo término de la progresión aritmética 20, 26, 32, \ldots es:
\item El vigésimo término de la progresión aritmética 18, 21, 24, \ldots es:
\item Número cuya dos cifras suman 8 y su diferencia es 2 y es primo.
\item El cuádruplo del quinto número primo.
\item El tercer término de una progresión geométrica cuya razón $r$ es 3 y cuyo primer término $a_{1}$ es 5.
\item El ángulo menor de un triángulo cuyos ángulos están en progresión aritmética cuya diferencia aritmética entre dos términos consecutivos es 4.

Seleccione la ecuación con la que se puede solucionar cada una de las siguientes situaciones.
\item El triple de un número aumentado en dos equivale a 3 unidades menos que el doble de dicho número. ¿Cuál es el número?
\begin{enumerate}
\item $x^{3}+2=3-x^{2}$ (58)
\item $3x+2=2x-3$ (66)
\item $3x+2x=2x-3$ (79)
\item $x^{3}+2=x^{2}-3$ (86)
\end{enumerate}
\item La suma de tres números impares consecutivos es 105. ¿Cuál es el número mayor?
\begin{enumerate}
\item $x+(x-2)+(x-4)=105$ (75)
\item $x+(x-1)+(x-2)=105$ (93)
\item $x+(x+1)+(x+2)=105$ (25)
\item $x+(x+2)+(x+4)=105$ (63)
\end{enumerate}
\item ¿Cuál número al sumarle 90 y dividirlo entre el mismo da como resultado 31?
\begin{enumerate}
\item $\dfrac{90+x}{90+x}=31$ (92)
\item $\dfrac{90x}{x}=31$ (87)
\item $\dfrac{90x}{90x}=31$ (54)
\item $\dfrac{x+90}{x}=31$ (77)
\end{enumerate}
\item El sexto término de la sucesión cuyo término general es $a_{n}=2^{n}+4$
\item El área del cuadrado ABCD es 3600 $u^{2}$. ¿El valor de $m$ es?
\begin{center}
\begin{tikzpicture}
\draw (0,0)node[left]{A}--(1.41,1.41)node[above]{B}--(2.82,0)node[right]{C}--(1.41,-1.41)node[below]{D}--cycle;
\draw (0,0)--node[left]{36u}(0,1.41)--node[above]{m}(1.41,1.41);
\draw (0,1.21)--(.21,1.21)--(0.21,1.41);
\end{tikzpicture} 
\end{center}
\item La figura es simétrica y su perímetro es 364 $u$. ¿El valor de $q$ es?
\begin{center}
%Uncomment next line if XeTeX is used
%\def\pgfsysdriver{pgfsys-xetex.def}
\usetikzlibrary{arrows}
\baselineskip=10pt
\hsize=6.3truein
\vsize=8.7truein
\definecolor{zzttqq}{rgb}{0.26666666666666666,0.26666666666666666,0.26666666666666666}
\tikzpicture[scale=.65,line cap=round,line join=round,>=triangle 45,x=1.0cm,y=1.0cm]
\clip(-4.6,-1.86) rectangle (6.94,6.54);
\fill[color=zzttqq,fill=zzttqq,fill opacity=0.1] (-1.,6.) -- (2.,6.) -- (2.,4.) -- (4.,4.) -- (4.,1.) -- (2.,1.) -- (2.,-1.) -- (-1.,-1.) -- (-1.,1.) -- (-3.,1.) -- (-3.,4.) -- (-1.,4.) -- cycle;
\draw [color=zzttqq] (-1.,6.)-- (2.,6.);
\draw [color=zzttqq] (2.,6.)-- (2.,4.);
\draw [color=zzttqq] (2.,4.)-- (4.,4.);
\draw [color=zzttqq] (4.,4.)-- (4.,1.);
\draw [color=zzttqq] (4.,1.)-- (2.,1.);
\draw [color=zzttqq] (2.,1.)-- (2.,-1.);
\draw [color=zzttqq] (2.,-1.)-- (-1.,-1.);
\draw [color=zzttqq] (-1.,-1.)-- (-1.,1.);
\draw [color=zzttqq] (-1.,1.)-- (-3.,1.);
\draw [color=zzttqq] (-3.,1.)-- (-3.,4.);
\draw [color=zzttqq] (-3.,4.)-- (-1.,4.);
\draw [color=zzttqq] (-1.,4.)-- (-1.,6.);
\draw (0.28,6.45) node[anchor=north west] {39};
\draw (4,2.78) node[anchor=north west] {39};
\draw (2,5.24) node[anchor=north west] {q};
\draw (2.84,4.5) node[anchor=north west] {q};
\endtikzpicture
\end{center}
\item En años, la cuarta parte de un siglo.

Para los puntos 30--38, complete las siguientes tablas sabiendo que el producto entre los números de cada columna se indica en la parte inferior y la suma de los números de cada fila se indica a la derecha de la tabla.
\begin{center}
% Graphic for TeX using PGF
% Title: /home/german/Diagrama1.dia
% Creator: Dia v0.97.2
% CreationDate: Thu Jul 10 18:54:59 2014
% For: german
% \usepackage{tikz}
% The following commands are not supported in PSTricks at present
% We define them conditionally, so when they are implemented,
% this pgf file will use them.
\ifx\du\undefined
  \newlength{\du}
\fi
\setlength{\du}{15\unitlength}
\begin{tikzpicture}
\pgftransformxscale{1.000000}
\pgftransformyscale{-1.000000}
\definecolor{dialinecolor}{rgb}{0.000000, 0.000000, 0.000000}
\pgfsetstrokecolor{dialinecolor}
\definecolor{dialinecolor}{rgb}{1.000000, 1.000000, 1.000000}
\pgfsetfillcolor{dialinecolor}
% setfont left to latex
\definecolor{dialinecolor}{rgb}{0.000000, 0.000000, 0.000000}
\pgfsetstrokecolor{dialinecolor}
\node[anchor=west] at (13.250000\du,5.750000\du){};
% setfont left to latex
\definecolor{dialinecolor}{rgb}{0.000000, 0.000000, 0.000000}
\pgfsetstrokecolor{dialinecolor}
\node[anchor=west] at (20.300000\du,5.250000\du){Presentan};
\pgfsetlinewidth{0.100000\du}
\pgfsetdash{}{0pt}
\pgfsetdash{}{0pt}
\pgfsetbuttcap
{
\definecolor{dialinecolor}{rgb}{0.000000, 0.000000, 0.000000}
\pgfsetfillcolor{dialinecolor}
% was here!!!
\definecolor{dialinecolor}{rgb}{0.000000, 0.000000, 0.000000}
\pgfsetstrokecolor{dialinecolor}
\draw (21.564569\du,3.749960\du)--(21.600000\du,4.712500\du);
}
\pgfsetlinewidth{0.100000\du}
\pgfsetdash{}{0pt}
\pgfsetdash{}{0pt}
\pgfsetbuttcap
{
\definecolor{dialinecolor}{rgb}{0.000000, 0.000000, 0.000000}
\pgfsetfillcolor{dialinecolor}
% was here!!!
\pgfsetarrowsend{to}
\definecolor{dialinecolor}{rgb}{0.000000, 0.000000, 0.000000}
\pgfsetstrokecolor{dialinecolor}
\draw (21.550000\du,5.512500\du)--(21.616010\du,6.799686\du);
}
\pgfsetlinewidth{0.100000\du}
\pgfsetdash{}{0pt}
\pgfsetdash{}{0pt}
\pgfsetbuttcap
\pgfsetmiterjoin
\pgfsetlinewidth{0.100000\du}
\pgfsetbuttcap
\pgfsetmiterjoin
\pgfsetdash{}{0pt}
\definecolor{dialinecolor}{rgb}{1.000000, 1.000000, 1.000000}
\pgfsetfillcolor{dialinecolor}
\pgfpathmoveto{\pgfpoint{17.268750\du}{6.687500\du}}
\pgfpathlineto{\pgfpoint{25.881250\du}{6.687500\du}}
\pgfpathcurveto{\pgfpoint{27.070389\du}{6.687500\du}}{\pgfpoint{28.034375\du}{7.193978\du}}{\pgfpoint{28.034375\du}{7.818750\du}}
\pgfpathcurveto{\pgfpoint{28.034375\du}{8.443522\du}}{\pgfpoint{27.070389\du}{8.950000\du}}{\pgfpoint{25.881250\du}{8.950000\du}}
\pgfpathlineto{\pgfpoint{17.268750\du}{8.950000\du}}
\pgfpathcurveto{\pgfpoint{16.079611\du}{8.950000\du}}{\pgfpoint{15.115625\du}{8.443522\du}}{\pgfpoint{15.115625\du}{7.818750\du}}
\pgfpathcurveto{\pgfpoint{15.115625\du}{7.193978\du}}{\pgfpoint{16.079611\du}{6.687500\du}}{\pgfpoint{17.268750\du}{6.687500\du}}
\pgfusepath{fill}
\definecolor{dialinecolor}{rgb}{0.000000, 0.000000, 0.000000}
\pgfsetstrokecolor{dialinecolor}
\pgfpathmoveto{\pgfpoint{17.268750\du}{6.687500\du}}
\pgfpathlineto{\pgfpoint{25.881250\du}{6.687500\du}}
\pgfpathcurveto{\pgfpoint{27.070389\du}{6.687500\du}}{\pgfpoint{28.034375\du}{7.193978\du}}{\pgfpoint{28.034375\du}{7.818750\du}}
\pgfpathcurveto{\pgfpoint{28.034375\du}{8.443522\du}}{\pgfpoint{27.070389\du}{8.950000\du}}{\pgfpoint{25.881250\du}{8.950000\du}}
\pgfpathlineto{\pgfpoint{17.268750\du}{8.950000\du}}
\pgfpathcurveto{\pgfpoint{16.079611\du}{8.950000\du}}{\pgfpoint{15.115625\du}{8.443522\du}}{\pgfpoint{15.115625\du}{7.818750\du}}
\pgfpathcurveto{\pgfpoint{15.115625\du}{7.193978\du}}{\pgfpoint{16.079611\du}{6.687500\du}}{\pgfpoint{17.268750\du}{6.687500\du}}
\pgfusepath{stroke}
% setfont left to latex
\definecolor{dialinecolor}{rgb}{0.000000, 0.000000, 0.000000}
\pgfsetstrokecolor{dialinecolor}
\node at (21.575000\du,8.018750\du){Relaciones multiplicativas};
% setfont left to latex
\definecolor{dialinecolor}{rgb}{0.000000, 0.000000, 0.000000}
\pgfsetstrokecolor{dialinecolor}
\node[anchor=west] at (21.575000\du,7.818750\du){};
\pgfsetlinewidth{0.100000\du}
\pgfsetdash{}{0pt}
\pgfsetdash{}{0pt}
\pgfsetbuttcap
\pgfsetmiterjoin
\pgfsetlinewidth{0.100000\du}
\pgfsetbuttcap
\pgfsetmiterjoin
\pgfsetdash{}{0pt}
\definecolor{dialinecolor}{rgb}{1.000000, 1.000000, 1.000000}
\pgfsetfillcolor{dialinecolor}
\pgfpathmoveto{\pgfpoint{16.998750\du}{1.650000\du}}
\pgfpathlineto{\pgfpoint{26.051250\du}{1.650000\du}}
\pgfpathcurveto{\pgfpoint{27.301140\du}{1.650000\du}}{\pgfpoint{28.314375\du}{2.108908\du}}{\pgfpoint{28.314375\du}{2.675000\du}}
\pgfpathcurveto{\pgfpoint{28.314375\du}{3.241092\du}}{\pgfpoint{27.301140\du}{3.700000\du}}{\pgfpoint{26.051250\du}{3.700000\du}}
\pgfpathlineto{\pgfpoint{16.998750\du}{3.700000\du}}
\pgfpathcurveto{\pgfpoint{15.748860\du}{3.700000\du}}{\pgfpoint{14.735625\du}{3.241092\du}}{\pgfpoint{14.735625\du}{2.675000\du}}
\pgfpathcurveto{\pgfpoint{14.735625\du}{2.108908\du}}{\pgfpoint{15.748860\du}{1.650000\du}}{\pgfpoint{16.998750\du}{1.650000\du}}
\pgfusepath{fill}
\definecolor{dialinecolor}{rgb}{0.000000, 0.000000, 0.000000}
\pgfsetstrokecolor{dialinecolor}
\pgfpathmoveto{\pgfpoint{16.998750\du}{1.650000\du}}
\pgfpathlineto{\pgfpoint{26.051250\du}{1.650000\du}}
\pgfpathcurveto{\pgfpoint{27.301140\du}{1.650000\du}}{\pgfpoint{28.314375\du}{2.108908\du}}{\pgfpoint{28.314375\du}{2.675000\du}}
\pgfpathcurveto{\pgfpoint{28.314375\du}{3.241092\du}}{\pgfpoint{27.301140\du}{3.700000\du}}{\pgfpoint{26.051250\du}{3.700000\du}}
\pgfpathlineto{\pgfpoint{16.998750\du}{3.700000\du}}
\pgfpathcurveto{\pgfpoint{15.748860\du}{3.700000\du}}{\pgfpoint{14.735625\du}{3.241092\du}}{\pgfpoint{14.735625\du}{2.675000\du}}
\pgfpathcurveto{\pgfpoint{14.735625\du}{2.108908\du}}{\pgfpoint{15.748860\du}{1.650000\du}}{\pgfpoint{16.998750\du}{1.650000\du}}
\pgfusepath{stroke}
% setfont left to latex
\definecolor{dialinecolor}{rgb}{0.000000, 0.000000, 0.000000}
\pgfsetstrokecolor{dialinecolor}
\node at (21.525000\du,2.914712\du){Los números naturales};
\pgfsetlinewidth{0.100000\du}
\pgfsetdash{}{0pt}
\pgfsetdash{}{0pt}
\pgfsetbuttcap
{
\definecolor{dialinecolor}{rgb}{0.000000, 0.000000, 0.000000}
\pgfsetfillcolor{dialinecolor}
% was here!!!
\definecolor{dialinecolor}{rgb}{0.000000, 0.000000, 0.000000}
\pgfsetstrokecolor{dialinecolor}
\draw (21.607263\du,8.999045\du)--(21.650000\du,10.562500\du);
}
% setfont left to latex
\definecolor{dialinecolor}{rgb}{0.000000, 0.000000, 0.000000}
\pgfsetstrokecolor{dialinecolor}
\node[anchor=west] at (20.800000\du,11.312500\du){como};
\pgfsetlinewidth{0.100000\du}
\pgfsetdash{}{0pt}
\pgfsetdash{}{0pt}
\pgfsetbuttcap
{
\definecolor{dialinecolor}{rgb}{0.000000, 0.000000, 0.000000}
\pgfsetfillcolor{dialinecolor}
% was here!!!
\pgfsetarrowsend{to}
\definecolor{dialinecolor}{rgb}{0.000000, 0.000000, 0.000000}
\pgfsetstrokecolor{dialinecolor}
\draw (21.650000\du,11.612500\du)--(21.625000\du,13.050000\du);
}
\pgfsetlinewidth{0.100000\du}
\pgfsetdash{}{0pt}
\pgfsetdash{}{0pt}
\pgfsetbuttcap
{
\definecolor{dialinecolor}{rgb}{0.000000, 0.000000, 0.000000}
\pgfsetfillcolor{dialinecolor}
% was here!!!
\definecolor{dialinecolor}{rgb}{0.000000, 0.000000, 0.000000}
\pgfsetstrokecolor{dialinecolor}
\draw (16.500000\du,13.050000\du)--(26.750000\du,13.050000\du);
}
\pgfsetlinewidth{0.100000\du}
\pgfsetdash{}{0pt}
\pgfsetdash{}{0pt}
\pgfsetbuttcap
{
\definecolor{dialinecolor}{rgb}{0.000000, 0.000000, 0.000000}
\pgfsetfillcolor{dialinecolor}
% was here!!!
\definecolor{dialinecolor}{rgb}{0.000000, 0.000000, 0.000000}
\pgfsetstrokecolor{dialinecolor}
\draw (16.600000\du,13.087500\du)--(16.549658\du,14.150970\du);
}
\pgfsetlinewidth{0.100000\du}
\pgfsetdash{}{0pt}
\pgfsetdash{}{0pt}
\pgfsetbuttcap
\pgfsetmiterjoin
\pgfsetlinewidth{0.100000\du}
\pgfsetbuttcap
\pgfsetmiterjoin
\pgfsetdash{}{0pt}
\definecolor{dialinecolor}{rgb}{1.000000, 1.000000, 1.000000}
\pgfsetfillcolor{dialinecolor}
\pgfpathmoveto{\pgfpoint{13.503750\du}{14.200000\du}}
\pgfpathlineto{\pgfpoint{19.496250\du}{14.200000\du}}
\pgfpathcurveto{\pgfpoint{20.323642\du}{14.200000\du}}{\pgfpoint{20.994375\du}{14.647715\du}}{\pgfpoint{20.994375\du}{15.200000\du}}
\pgfpathcurveto{\pgfpoint{20.994375\du}{15.752285\du}}{\pgfpoint{20.323642\du}{16.200000\du}}{\pgfpoint{19.496250\du}{16.200000\du}}
\pgfpathlineto{\pgfpoint{13.503750\du}{16.200000\du}}
\pgfpathcurveto{\pgfpoint{12.676358\du}{16.200000\du}}{\pgfpoint{12.005625\du}{15.752285\du}}{\pgfpoint{12.005625\du}{15.200000\du}}
\pgfpathcurveto{\pgfpoint{12.005625\du}{14.647715\du}}{\pgfpoint{12.676358\du}{14.200000\du}}{\pgfpoint{13.503750\du}{14.200000\du}}
\pgfusepath{fill}
\definecolor{dialinecolor}{rgb}{0.000000, 0.000000, 0.000000}
\pgfsetstrokecolor{dialinecolor}
\pgfpathmoveto{\pgfpoint{13.503750\du}{14.200000\du}}
\pgfpathlineto{\pgfpoint{19.496250\du}{14.200000\du}}
\pgfpathcurveto{\pgfpoint{20.323642\du}{14.200000\du}}{\pgfpoint{20.994375\du}{14.647715\du}}{\pgfpoint{20.994375\du}{15.200000\du}}
\pgfpathcurveto{\pgfpoint{20.994375\du}{15.752285\du}}{\pgfpoint{20.323642\du}{16.200000\du}}{\pgfpoint{19.496250\du}{16.200000\du}}
\pgfpathlineto{\pgfpoint{13.503750\du}{16.200000\du}}
\pgfpathcurveto{\pgfpoint{12.676358\du}{16.200000\du}}{\pgfpoint{12.005625\du}{15.752285\du}}{\pgfpoint{12.005625\du}{15.200000\du}}
\pgfpathcurveto{\pgfpoint{12.005625\du}{14.647715\du}}{\pgfpoint{12.676358\du}{14.200000\du}}{\pgfpoint{13.503750\du}{14.200000\du}}
\pgfusepath{stroke}
% setfont left to latex
\definecolor{dialinecolor}{rgb}{0.000000, 0.000000, 0.000000}
\pgfsetstrokecolor{dialinecolor}
\node at (16.500000\du,15.400000\du){Ser múltiplo de ...};
\pgfsetlinewidth{0.100000\du}
\pgfsetdash{}{0pt}
\pgfsetdash{}{0pt}
\pgfsetbuttcap
\pgfsetmiterjoin
\pgfsetlinewidth{0.100000\du}
\pgfsetbuttcap
\pgfsetmiterjoin
\pgfsetdash{}{0pt}
\definecolor{dialinecolor}{rgb}{1.000000, 1.000000, 1.000000}
\pgfsetfillcolor{dialinecolor}
\pgfpathmoveto{\pgfpoint{23.943750\du}{14.300000\du}}
\pgfpathlineto{\pgfpoint{29.456250\du}{14.300000\du}}
\pgfpathcurveto{\pgfpoint{30.217368\du}{14.300000\du}}{\pgfpoint{30.834375\du}{14.736522\du}}{\pgfpoint{30.834375\du}{15.275000\du}}
\pgfpathcurveto{\pgfpoint{30.834375\du}{15.813478\du}}{\pgfpoint{30.217368\du}{16.250000\du}}{\pgfpoint{29.456250\du}{16.250000\du}}
\pgfpathlineto{\pgfpoint{23.943750\du}{16.250000\du}}
\pgfpathcurveto{\pgfpoint{23.182632\du}{16.250000\du}}{\pgfpoint{22.565625\du}{15.813478\du}}{\pgfpoint{22.565625\du}{15.275000\du}}
\pgfpathcurveto{\pgfpoint{22.565625\du}{14.736522\du}}{\pgfpoint{23.182632\du}{14.300000\du}}{\pgfpoint{23.943750\du}{14.300000\du}}
\pgfusepath{fill}
\definecolor{dialinecolor}{rgb}{0.000000, 0.000000, 0.000000}
\pgfsetstrokecolor{dialinecolor}
\pgfpathmoveto{\pgfpoint{23.943750\du}{14.300000\du}}
\pgfpathlineto{\pgfpoint{29.456250\du}{14.300000\du}}
\pgfpathcurveto{\pgfpoint{30.217368\du}{14.300000\du}}{\pgfpoint{30.834375\du}{14.736522\du}}{\pgfpoint{30.834375\du}{15.275000\du}}
\pgfpathcurveto{\pgfpoint{30.834375\du}{15.813478\du}}{\pgfpoint{30.217368\du}{16.250000\du}}{\pgfpoint{29.456250\du}{16.250000\du}}
\pgfpathlineto{\pgfpoint{23.943750\du}{16.250000\du}}
\pgfpathcurveto{\pgfpoint{23.182632\du}{16.250000\du}}{\pgfpoint{22.565625\du}{15.813478\du}}{\pgfpoint{22.565625\du}{15.275000\du}}
\pgfpathcurveto{\pgfpoint{22.565625\du}{14.736522\du}}{\pgfpoint{23.182632\du}{14.300000\du}}{\pgfpoint{23.943750\du}{14.300000\du}}
\pgfusepath{stroke}
% setfont left to latex
\definecolor{dialinecolor}{rgb}{0.000000, 0.000000, 0.000000}
\pgfsetstrokecolor{dialinecolor}
\node at (26.700000\du,15.475000\du){Ser divisor de ...};
\pgfsetlinewidth{0.100000\du}
\pgfsetdash{}{0pt}
\pgfsetdash{}{0pt}
\pgfsetbuttcap
{
\definecolor{dialinecolor}{rgb}{0.000000, 0.000000, 0.000000}
\pgfsetfillcolor{dialinecolor}
% was here!!!
\definecolor{dialinecolor}{rgb}{0.000000, 0.000000, 0.000000}
\pgfsetstrokecolor{dialinecolor}
\draw (26.700000\du,12.987500\du)--(26.700000\du,14.250763\du);
}
\pgfsetlinewidth{0.100000\du}
\pgfsetdash{}{0pt}
\pgfsetdash{}{0pt}
\pgfsetbuttcap
{
\definecolor{dialinecolor}{rgb}{0.000000, 0.000000, 0.000000}
\pgfsetfillcolor{dialinecolor}
% was here!!!
\definecolor{dialinecolor}{rgb}{0.000000, 0.000000, 0.000000}
\pgfsetstrokecolor{dialinecolor}
\draw (26.700000\du,16.300089\du)--(26.700000\du,17.937500\du);
}
% setfont left to latex
\definecolor{dialinecolor}{rgb}{0.000000, 0.000000, 0.000000}
\pgfsetstrokecolor{dialinecolor}
\node[anchor=west] at (25.100000\du,18.537500\du){determina};
\pgfsetlinewidth{0.100000\du}
\pgfsetdash{}{0pt}
\pgfsetdash{}{0pt}
\pgfsetbuttcap
{
\definecolor{dialinecolor}{rgb}{0.000000, 0.000000, 0.000000}
\pgfsetfillcolor{dialinecolor}
% was here!!!
\pgfsetarrowsend{to}
\definecolor{dialinecolor}{rgb}{0.000000, 0.000000, 0.000000}
\pgfsetstrokecolor{dialinecolor}
\draw (26.700000\du,18.687500\du)--(26.775000\du,20.412500\du);
}
\pgfsetlinewidth{0.100000\du}
\pgfsetdash{}{0pt}
\pgfsetdash{}{0pt}
\pgfsetbuttcap
{
\definecolor{dialinecolor}{rgb}{0.000000, 0.000000, 0.000000}
\pgfsetfillcolor{dialinecolor}
% was here!!!
\definecolor{dialinecolor}{rgb}{0.000000, 0.000000, 0.000000}
\pgfsetstrokecolor{dialinecolor}
\draw (21.750000\du,20.437500\du)--(31.800000\du,20.387500\du);
}
\pgfsetlinewidth{0.100000\du}
\pgfsetdash{}{0pt}
\pgfsetdash{}{0pt}
\pgfsetbuttcap
{
\definecolor{dialinecolor}{rgb}{0.000000, 0.000000, 0.000000}
\pgfsetfillcolor{dialinecolor}
% was here!!!
\definecolor{dialinecolor}{rgb}{0.000000, 0.000000, 0.000000}
\pgfsetstrokecolor{dialinecolor}
\draw (21.850000\du,20.487500\du)--(21.876843\du,21.437080\du);
}
\pgfsetlinewidth{0.100000\du}
\pgfsetdash{}{0pt}
\pgfsetdash{}{0pt}
\pgfsetbuttcap
\pgfsetmiterjoin
\pgfsetlinewidth{0.100000\du}
\pgfsetbuttcap
\pgfsetmiterjoin
\pgfsetdash{}{0pt}
\definecolor{dialinecolor}{rgb}{1.000000, 1.000000, 1.000000}
\pgfsetfillcolor{dialinecolor}
\pgfpathmoveto{\pgfpoint{19.113750\du}{21.487500\du}}
\pgfpathlineto{\pgfpoint{24.686250\du}{21.487500\du}}
\pgfpathcurveto{\pgfpoint{25.455652\du}{21.487500\du}}{\pgfpoint{26.079375\du}{21.831681\du}}{\pgfpoint{26.079375\du}{22.256250\du}}
\pgfpathcurveto{\pgfpoint{26.079375\du}{22.680819\du}}{\pgfpoint{25.455652\du}{23.025000\du}}{\pgfpoint{24.686250\du}{23.025000\du}}
\pgfpathlineto{\pgfpoint{19.113750\du}{23.025000\du}}
\pgfpathcurveto{\pgfpoint{18.344348\du}{23.025000\du}}{\pgfpoint{17.720625\du}{22.680819\du}}{\pgfpoint{17.720625\du}{22.256250\du}}
\pgfpathcurveto{\pgfpoint{17.720625\du}{21.831681\du}}{\pgfpoint{18.344348\du}{21.487500\du}}{\pgfpoint{19.113750\du}{21.487500\du}}
\pgfusepath{fill}
\definecolor{dialinecolor}{rgb}{0.000000, 0.000000, 0.000000}
\pgfsetstrokecolor{dialinecolor}
\pgfpathmoveto{\pgfpoint{19.113750\du}{21.487500\du}}
\pgfpathlineto{\pgfpoint{24.686250\du}{21.487500\du}}
\pgfpathcurveto{\pgfpoint{25.455652\du}{21.487500\du}}{\pgfpoint{26.079375\du}{21.831681\du}}{\pgfpoint{26.079375\du}{22.256250\du}}
\pgfpathcurveto{\pgfpoint{26.079375\du}{22.680819\du}}{\pgfpoint{25.455652\du}{23.025000\du}}{\pgfpoint{24.686250\du}{23.025000\du}}
\pgfpathlineto{\pgfpoint{19.113750\du}{23.025000\du}}
\pgfpathcurveto{\pgfpoint{18.344348\du}{23.025000\du}}{\pgfpoint{17.720625\du}{22.680819\du}}{\pgfpoint{17.720625\du}{22.256250\du}}
\pgfpathcurveto{\pgfpoint{17.720625\du}{21.831681\du}}{\pgfpoint{18.344348\du}{21.487500\du}}{\pgfpoint{19.113750\du}{21.487500\du}}
\pgfusepath{stroke}
% setfont left to latex
\definecolor{dialinecolor}{rgb}{0.000000, 0.000000, 0.000000}
\pgfsetstrokecolor{dialinecolor}
\node at (21.900000\du,22.456250\du){Números primos};
\pgfsetlinewidth{0.100000\du}
\pgfsetdash{}{0pt}
\pgfsetdash{}{0pt}
\pgfsetbuttcap
\pgfsetmiterjoin
\pgfsetlinewidth{0.100000\du}
\pgfsetbuttcap
\pgfsetmiterjoin
\pgfsetdash{}{0pt}
\definecolor{dialinecolor}{rgb}{1.000000, 1.000000, 1.000000}
\pgfsetfillcolor{dialinecolor}
\pgfpathmoveto{\pgfpoint{28.320000\du}{21.487500\du}}
\pgfpathlineto{\pgfpoint{35.580000\du}{21.487500\du}}
\pgfpathcurveto{\pgfpoint{36.582397\du}{21.487500\du}}{\pgfpoint{37.395000\du}{21.868058\du}}{\pgfpoint{37.395000\du}{22.337500\du}}
\pgfpathcurveto{\pgfpoint{37.395000\du}{22.806942\du}}{\pgfpoint{36.582397\du}{23.187500\du}}{\pgfpoint{35.580000\du}{23.187500\du}}
\pgfpathlineto{\pgfpoint{28.320000\du}{23.187500\du}}
\pgfpathcurveto{\pgfpoint{27.317603\du}{23.187500\du}}{\pgfpoint{26.505000\du}{22.806942\du}}{\pgfpoint{26.505000\du}{22.337500\du}}
\pgfpathcurveto{\pgfpoint{26.505000\du}{21.868058\du}}{\pgfpoint{27.317603\du}{21.487500\du}}{\pgfpoint{28.320000\du}{21.487500\du}}
\pgfusepath{fill}
\definecolor{dialinecolor}{rgb}{0.000000, 0.000000, 0.000000}
\pgfsetstrokecolor{dialinecolor}
\pgfpathmoveto{\pgfpoint{28.320000\du}{21.487500\du}}
\pgfpathlineto{\pgfpoint{35.580000\du}{21.487500\du}}
\pgfpathcurveto{\pgfpoint{36.582397\du}{21.487500\du}}{\pgfpoint{37.395000\du}{21.868058\du}}{\pgfpoint{37.395000\du}{22.337500\du}}
\pgfpathcurveto{\pgfpoint{37.395000\du}{22.806942\du}}{\pgfpoint{36.582397\du}{23.187500\du}}{\pgfpoint{35.580000\du}{23.187500\du}}
\pgfpathlineto{\pgfpoint{28.320000\du}{23.187500\du}}
\pgfpathcurveto{\pgfpoint{27.317603\du}{23.187500\du}}{\pgfpoint{26.505000\du}{22.806942\du}}{\pgfpoint{26.505000\du}{22.337500\du}}
\pgfpathcurveto{\pgfpoint{26.505000\du}{21.868058\du}}{\pgfpoint{27.317603\du}{21.487500\du}}{\pgfpoint{28.320000\du}{21.487500\du}}
\pgfusepath{stroke}
% setfont left to latex
\definecolor{dialinecolor}{rgb}{0.000000, 0.000000, 0.000000}
\pgfsetstrokecolor{dialinecolor}
\node at (31.950000\du,22.537500\du){Números compuestos};
\pgfsetlinewidth{0.100000\du}
\pgfsetdash{}{0pt}
\pgfsetdash{}{0pt}
\pgfsetbuttcap
{
\definecolor{dialinecolor}{rgb}{0.000000, 0.000000, 0.000000}
\pgfsetfillcolor{dialinecolor}
% was here!!!
\definecolor{dialinecolor}{rgb}{0.000000, 0.000000, 0.000000}
\pgfsetstrokecolor{dialinecolor}
\draw (31.850000\du,20.487500\du)--(31.901343\du,21.437341\du);
}
\end{tikzpicture}

\end{center}
\end{enumerate}
\end{multicols}
\newpage
\begin{tikzpicture}
 \fill (1,0) node[above]{1} circle (0.2ex);
 \fill (2,0) node[above]{2} circle (0.2ex);
 \fill (3,0) node[above]{3} circle (0.2ex);
 \fill (4,0) node[above]{4} circle (0.2ex);
 \fill (5,0) node[above]{5} circle (0.2ex);
 \fill (6,0) node[above]{6} circle (0.2ex);
 \fill (7,0) node[above]{7} circle (0.2ex);
 \fill (8,0) node[above]{8} circle (0.2ex);
 \fill (9,0) node[above]{9} circle (0.2ex);
 \fill (10,0) node[above]{10} circle (0.2ex);
 \fill (1,-1) node[left]{11} circle (0.2ex);
 \fill (2,-1) circle (0.2ex);
 \fill (3,-1) circle (0.2ex);
 \fill (4,-1) circle (0.2ex);
 \fill (5,-1) circle (0.2ex);
 \fill (6,-1) circle (0.2ex);
 \fill (7,-1) circle (0.2ex);
 \fill (8,-1) circle (0.2ex);
 \fill (9,-1) circle (0.2ex);
 \fill (10,-1) circle (0.2ex);
 \fill (1,-2) node[left]{21} circle (0.2ex);
 \fill (2,-2) circle (0.2ex);
 \fill (3,-2) circle (0.2ex);
 \fill (4,-2) circle (0.2ex);
 \fill (5,-2) circle (0.2ex);
 \fill (6,-2) circle (0.2ex);
 \fill (7,-2) circle (0.2ex);
 \fill (8,-2) circle (0.2ex);
 \fill (9,-2) circle (0.2ex);
 \fill (10,-2) circle (0.2ex);
 \fill (1,-3) node[left]{31} circle (0.2ex);
 \fill (2,-3) circle (0.2ex);
 \fill (3,-3) circle (0.2ex);
 \fill (4,-3) circle (0.2ex);
 \fill (5,-3) circle (0.2ex);
 \fill (6,-3) circle (0.2ex);
 \fill (7,-3) circle (0.2ex);
 \fill (8,-3) circle (0.2ex);
 \fill (9,-3) circle (0.2ex);
 \fill (10,-3) circle (0.2ex);
 \fill (1,-4) node[left]{41} circle (0.2ex);
 \fill (2,-4) circle (0.2ex);
 \fill (3,-4) circle (0.2ex);
 \fill (4,-4) circle (0.2ex);
 \fill (5,-4) circle (0.2ex);
 \fill (6,-4) circle (0.2ex);
 \fill (7,-4) circle (0.2ex);
 \fill (8,-4) circle (0.2ex);
 \fill (9,-4) circle (0.2ex);
 \fill (10,-4) node[right]{50} circle (0.2ex);
 \fill (1,-5) node[left]{51} circle (0.2ex);
 \fill (2,-5) circle (0.2ex);
 \fill (3,-5) circle (0.2ex);
 \fill (4,-5) circle (0.2ex);
 \fill (5,-5) circle (0.2ex);
 \fill (6,-5) circle (0.2ex);
 \fill (7,-5) circle (0.2ex);
 \fill (8,-5) circle (0.2ex);
 \fill (9,-5) circle (0.2ex);
 \fill (10,-5) circle (0.2ex);
 \fill (1,-6) node[left]{61} circle (0.2ex);
 \fill (2,-6) circle (0.2ex);
 \fill (3,-6) circle (0.2ex);
 \fill (4,-6) circle (0.2ex);
 \fill (5,-6) circle (0.2ex);
 \fill (6,-6) circle (0.2ex);
 \fill (7,-6) circle (0.2ex);
 \fill (8,-6) circle (0.2ex);
 \fill (9,-6) circle (0.2ex);
 \fill (10,-6) circle (0.2ex);
 \fill (1,-7) node[left]{71} circle (0.2ex);
 \fill (2,-7) circle (0.2ex);
 \fill (3,-7) circle (0.2ex);
 \fill (4,-7) circle (0.2ex);
 \fill (5,-7) circle (0.2ex);
 \fill (6,-7) circle (0.2ex);
 \fill (7,-7) circle (0.2ex);
 \fill (8,-7) circle (0.2ex);
 \fill (9,-7) circle (0.2ex);
 \fill (10,-7) circle (0.2ex);
 \fill (1,-8) node[left]{81} circle (0.2ex);
 \fill (2,-8) circle (0.2ex);
 \fill (3,-8) circle (0.2ex);
 \fill (4,-8) circle (0.2ex);
 \fill (5,-8) circle (0.2ex);
 \fill (6,-8) circle (0.2ex);
 \fill (7,-8) circle (0.2ex);
 \fill (8,-8) circle (0.2ex);
 \fill (9,-8) circle (0.2ex);
 \fill (10,-8) circle (0.2ex);
 \fill (1,-9) node[left]{91} circle (0.2ex);
 \fill (2,-9) circle (0.2ex);
 \fill (3,-9) circle (0.2ex);
 \fill (4,-9) circle (0.2ex);
 \fill (5,-9) circle (0.2ex);
 \fill (6,-9) circle (0.2ex);
 \fill (7,-9) circle (0.2ex);
 \fill (8,-9) circle (0.2ex);
 \fill (9,-9) circle (0.2ex);
 \fill (10,-9) node[right]{100} circle (0.2ex);
 % \draw (2,-7)--(2,-8)--(3,-9)--(8,-9)--(9,-8)--(9,-3)--(8,-4)--(9,-4)--(10,-3)--(8,-3)--(7,-1)--(3,-1)--(2,-3)--(1,-4)--(2,-7)--(3,-6)--(4,-7)--(5,-7)--(3,-5)--(4,-4)--(5,-4)--(6,-5)--(6,-6)--(5,-7)--(7,-7)--(8,-6)--(8,-4)--(6,-2)--(5,-2)--(5,-1)--(6,0)--(6,-1)--(5,-2)--(3,-3)--(3,-7)--(5,-8)--(6,-7)--(7,-9);
\end{tikzpicture}

\end{document}
